% GENERATED FILE. Edit canonical lessons, then run npm run generate:books.
% canonical-source-hash: fnv1a64:37a9e4035ae1ff9d
% canonical-lessons: HI-C486-dukhna, HI-C486-pindli, HI-C486-kasakna, HI-C486-rirh, HI-C486-sujna

\chapter{Aches and Pains}
\label{ch:hi-aches-and-pains}
\hlchaptermodality{486}

\begin{chapteropening}
\textbf{By the end of this chapter:} \emph{I can say where it hurts and how.}

Complete the last lesson of chapter 486: I can say where it hurts and how.
\end{chapteropening}

\section[\texorpdfstring{\hi{दुखना} (dukhnā)}{दुखना (dukhnā)}]{\hi{दुखना} (dukhnā) — to ache}

\label{lesson:HI-C486-dukhna}



\begin{quote}
Before the new one: say the Hindi for to retch, then the Hindi for a burp.
\end{quote}

\subsection*{You'll want to know: \hi{दुखना}}
\textbf{\hi{दुखना}} — \emph{dukhnā} — \textquotedblleft{}to ache\textquotedblright{}.

\begin{quote}
\hi{ठंड} \hi{में} \hi{दादा} \hi{के} \hi{घुटने} \hi{दुखते} \hi{हैं।}
\end{quote}

\begin{quote}
\emph{ṭhaṇḍ meṁ dādā ke ghuṭne dukhte haiṁ.}
\end{quote}

\begin{quote}
Grandfather's knees ache in the cold.
\end{quote}

Past (it …): \hi{वह} \hi{पैर} \hi{बहुत} \hi{दुखा}, \emph{vah pair bahut dukhā}.

\begin{cousinweb}
\textbf{\hi{दुखना}} (\emph{dukhnā}) is a tadbhava: CDIAL lists Hindi \emph{dukhnā}, \textquotedblleft{}pains\textquotedblright{}, under Sanskrit \emph{duḥkhati} (CDIAL 6376), the entry that also holds \emph{dukhānā}. The part of the body that aches is the subject.
\end{cousinweb}

\begin{grammarlens}[title={What you've built}]
One more word for describing what you lived and felt.
\end{grammarlens}

\subsection*{Your turn}
\begin{itemize}
\raggedright
  \item \emph{Say it:} \emph{dukhnā}
  \item \emph{Say it:} \emph{dukhnā}, once more
  \item \emph{Say it:} \emph{ṭhaṇḍ meṁ dādā ke ghuṭne dukhte haiṁ.} — Grandfather's knees ache in the cold.
  \item \emph{Recall:} say the Hindi for to retch, then the Hindi for a burp, then say \emph{dukhnā} again
\end{itemize}

\begin{tcolorbox}[breakable,colback=teal!4,colframe=teal!35!black,title={Before you move on}]
What does \hi{दुखना} mean? (\textquotedblleft{}To ache\textquotedblright{}.) Say it once more.
\end{tcolorbox}
\section[\texorpdfstring{\hi{पिंडली} (piṇḍlī)}{पिंडली (piṇḍlī)}]{\hi{पिंडली} (piṇḍlī) — the calf of the leg}

\label{lesson:HI-C486-pindli}



\begin{quote}
Before the new one: say the Hindi for a burp, then the Hindi for to ache.
\end{quote}

\subsection*{You'll want to know: \hi{पिंडली}}
\textbf{\hi{पिंडली}} — \emph{piṇḍlī} — \textquotedblleft{}the calf of the leg\textquotedblright{}.

\begin{quote}
\hi{दौड़ने} \hi{के} \hi{बाद} \hi{अरुण} \hi{की} \hi{पिंडली} \hi{में} \hi{दर्द} \hi{है।}
\end{quote}

\begin{quote}
\emph{dauṛne ke bād Aruṇ kī piṇḍlī meṁ dard hai.}
\end{quote}

\begin{quote}
Arun's calf aches after running.
\end{quote}

Feminine.

\begin{cousinweb}
\textbf{\hi{पिंडली}} (\emph{piṇḍlī}) is a tadbhava: CDIAL lists Hindi \emph{piṇḍlī}, \textquotedblleft{}calf of leg, shin\textquotedblright{}, beside \emph{piṇḍā}, \textquotedblleft{}lump\textquotedblright{}, under Sanskrit \emph{píṇḍa} (CDIAL 8168): the rounded lump of muscle at the back of the leg.
\end{cousinweb}

\begin{grammarlens}[title={What you've built}]
One more word for describing what you lived and felt.
\end{grammarlens}

\subsection*{Your turn}
\begin{itemize}
\raggedright
  \item \emph{Say it:} \emph{piṇḍlī}
  \item \emph{Say it:} \emph{piṇḍlī}, once more
  \item \emph{Say it:} \emph{dauṛne ke bād Aruṇ kī piṇḍlī meṁ dard hai.} — Arun's calf aches after running.
  \item \emph{Recall:} say the Hindi for a burp, then the Hindi for to ache, then say \emph{piṇḍlī} again
\end{itemize}

\begin{tcolorbox}[breakable,colback=teal!4,colframe=teal!35!black,title={Before you move on}]
What does \hi{पिंडली} mean? (\textquotedblleft{}The calf of the leg\textquotedblright{}.) Say it once more.
\end{tcolorbox}
\section[\texorpdfstring{\hi{कसकना} (kasaknā)}{कसकना (kasaknā)}]{\hi{कसकना} (kasaknā) — to throb with pain}

\label{lesson:HI-C486-kasakna}



\begin{quote}
Before the new one: say the Hindi for to ache, then the Hindi for the calf of the leg.
\end{quote}

\subsection*{You'll want to know: \hi{कसकना}}
\textbf{\hi{कसकना}} — \emph{kasaknā} — \textquotedblleft{}to throb with pain\textquotedblright{}.

\begin{quote}
\hi{ठंड} \hi{में} \hi{पुरानी} \hi{चोट} \hi{कसकती} \hi{है।}
\end{quote}

\begin{quote}
\emph{ṭhaṇḍ meṁ purānī coṭ kasaktī hai.}
\end{quote}

\begin{quote}
An old injury throbs in the cold.
\end{quote}

Past (it …): \hi{वह} \hi{चोट} \hi{रात} \hi{भर} \hi{कसकी}, \emph{vah coṭ rāt bhar kaskī}.

\begin{cousinweb}
\textbf{\hi{कसकना}} (\emph{kasaknā}) is inherited: CDIAL lists Hindi \emph{kasaknā}, \textquotedblleft{}to be racked by pain, ache\textquotedblright{}, under Sanskrit \emph{kaṣati}, \textquotedblleft{}rubs, scratches\textquotedblright{} (CDIAL 2972). It is the dull, returning ache of an old hurt.
\end{cousinweb}

\begin{grammarlens}[title={What you've built}]
One more word for describing what you lived and felt.
\end{grammarlens}

\subsection*{Your turn}
\begin{itemize}
\raggedright
  \item \emph{Say it:} \emph{kasaknā}
  \item \emph{Say it:} \emph{kasaknā}, once more
  \item \emph{Say it:} \emph{ṭhaṇḍ meṁ purānī coṭ kasaktī hai.} — An old injury throbs in the cold.
  \item \emph{Recall:} say the Hindi for to ache, then the Hindi for the calf of the leg, then say \emph{kasaknā} again
\end{itemize}

\begin{tcolorbox}[breakable,colback=teal!4,colframe=teal!35!black,title={Before you move on}]
What does \hi{कसकना} mean? (\textquotedblleft{}To throb with pain\textquotedblright{}.) Say it once more.
\end{tcolorbox}
\section[\texorpdfstring{\hi{रीढ़} (rīṛh)}{रीढ़ (rīṛh)}]{\hi{रीढ़} (rīṛh) — the spine; the backbone}

\label{lesson:HI-C486-rirh}



\begin{quote}
Before the new one: say the Hindi for the calf of the leg, then the Hindi for to throb with pain.
\end{quote}

\subsection*{You'll want to know: \hi{रीढ़}}
\textbf{\hi{रीढ़}} — \emph{rīṛh} — \textquotedblleft{}the spine; the backbone\textquotedblright{}.

\begin{quote}
\hi{भारी} \hi{बोझ} \hi{उठाने} \hi{से} \hi{मज़दूर} \hi{की} \hi{रीढ़} \hi{में} \hi{दर्द} \hi{है।}
\end{quote}

\begin{quote}
\emph{bhārī bojh uṭhāne se mazdūr kī rīṛh meṁ dard hai.}
\end{quote}

\begin{quote}
The labourer's back hurts from lifting heavy loads.
\end{quote}

Feminine.

\begin{cousinweb}
\textbf{\hi{रीढ़}} (\emph{rīṛh}) is a tadbhava: CDIAL lists Hindi \emph{rīṛh}, \textquotedblleft{}backbone\textquotedblright{}, under Sanskrit \emph{rīḍhaka-} (CDIAL 10749a). \emph{rīṛh kī haḍḍī}, \textquotedblleft{}the bone of the back\textquotedblright{}, is the spine itself.
\end{cousinweb}

\begin{grammarlens}[title={What you've built}]
One more word for describing what you lived and felt.
\end{grammarlens}

\subsection*{Your turn}
\begin{itemize}
\raggedright
  \item \emph{Say it:} \emph{rīṛh}
  \item \emph{Say it:} \emph{rīṛh}, once more
  \item \emph{Say it:} \emph{bhārī bojh uṭhāne se mazdūr kī rīṛh meṁ dard hai.} — The labourer's back hurts from lifting heavy loads.
  \item \emph{Recall:} say the Hindi for the calf of the leg, then the Hindi for to throb with pain, then say \emph{rīṛh} again
\end{itemize}

\begin{tcolorbox}[breakable,colback=teal!4,colframe=teal!35!black,title={Before you move on}]
What does \hi{रीढ़} mean? (\textquotedblleft{}The spine; the backbone\textquotedblright{}.) Say it once more.
\end{tcolorbox}
\section[\texorpdfstring{\hi{सूजना} (sūjnā)}{सूजना (sūjnā)}]{\hi{सूजना} (sūjnā) — to swell (of the body)}

\label{lesson:HI-C486-sujna}



\begin{quote}
Before the new one: say the Hindi for to throb with pain, then the Hindi for the spine.
\end{quote}

\subsection*{You'll want to know: \hi{सूजना}}
\textbf{\hi{सूजना}} — \emph{sūjnā} — \textquotedblleft{}to swell (of the body)\textquotedblright{}.

\begin{quote}
\hi{मोच} \hi{से} \hi{अरुण} \hi{का} \hi{पैर} \hi{सूज} \hi{गया।}
\end{quote}

\begin{quote}
\emph{moc se Aruṇ kā pair sūj gayā.}
\end{quote}

\begin{quote}
Arun's foot swelled up from the sprain.
\end{quote}

Past (it …): \hi{वह} \hi{पैर} \hi{सूज} \hi{गया}, \emph{vah pair sūj gayā}.

\begin{cousinweb}
\textbf{\hi{सूजना}} (\emph{sūjnā}) is a tadbhava: CDIAL lists Hindi \emph{sūjnā}, \textquotedblleft{}to swell\textquotedblright{}, under Sanskrit \emph{śūyatē}, \textquotedblleft{}swells\textquotedblright{} (CDIAL 12568). \emph{sūjan}, \textquotedblleft{}a swelling\textquotedblright{}, is its noun.
\end{cousinweb}

\begin{grammarlens}[title={What you've built}]
One more word for describing what you lived and felt.
\end{grammarlens}

\subsection*{Your turn}
\begin{itemize}
\raggedright
  \item \emph{Say it:} \emph{sūjnā}
  \item \emph{Say it:} \emph{sūjnā}, once more
  \item \emph{Say it:} \emph{moc se Aruṇ kā pair sūj gayā.} — Arun's foot swelled up from the sprain.
  \item \emph{Recall:} say the Hindi for to throb with pain, then the Hindi for the spine, then say \emph{sūjnā} again
\end{itemize}

\begin{tcolorbox}[breakable,colback=teal!4,colframe=teal!35!black,title={Before you move on}]
What does \hi{सूजना} mean? (\textquotedblleft{}To swell (of the body)\textquotedblright{}.) Say it once more.
\end{tcolorbox}
